% !TeX root = ../navier-stokes-manuscript.tex
A smooth fluid can only fail globally by reaching a last smooth time.
The problem is to understand whether the same three-dimensional velocity field can keep concentrating fast enough to create such an endpoint while viscosity is already diffusing the gradients it produces.
That field and its pressure satisfy
\begin{equation}\label{eq:NavierStokesSystem}
  \UseReadableDisplayRows
  \begin{gathered}
    \partial_tu+(u\cdot\nabla)u+\nabla p=\nu\Delta u,
    \qquad
    \nabla\cdot u=0
    \quad\text{on }\R^3\times(0,\infty),
    \\
    u(\cdot,0)=u_0
    \quad\text{on }\R^3.
  \end{gathered}
\end{equation}
Here \(u\) is the velocity, \(p\) is the pressure required by incompressibility, and the viscosity \(\nu>0\) stays fixed throughout the argument.
\MajorThoughtBreak
Fefferman's whole-space problem begins with a smooth divergence-free velocity whose spatial derivatives decay faster than every algebraic power \parencite{Fefferman2000}.
Write this datum class as
\begin{equation}\label{eq:FeffermanDatumClass}
  \UseReadableDisplayRows
  \begin{gathered}
    \Ssigma
    :=
    \mathscr S(\R^3;\R^3)\cap\ker(\nabla\cdot),
    \qquad
    u_0\in\Ssigma,
    \\
    \forall\alpha\in\Nzero^3\;
    \forall K\in\Nzero\;
    \exists C_{\alpha,K}<\infty\;
    \forall x\in\R^3,
    \qquad
    |\partial_x^\alpha u_0(x)|
    \leq
    C_{\alpha,K}(1+|x|)^{-K}.
  \end{gathered}
\end{equation}
The multi-index \(\alpha\) chooses a spatial derivative, and \(K\) chooses how much algebraic decay to demand.
For the fixed datum, \(C_{\alpha,K}\) may depend on those choices but never on \(x\).
\MajorThoughtBreak
For finite \(m\in\Nzero\), write
\[
  H^m_\sigma(\R^3)
  :=
  H^m(\R^3;\R^3)\cap\ker(\nabla\cdot).
\]
This rapid decay already places the datum at every finite Sobolev height.
Indeed, fix \(K>3/2\).  For every \(|\alpha|\leq m\), \eqref{eq:FeffermanDatumClass} gives
\[
  \norm{\partial_x^\alpha u_0}_2^2
  \leq
  C_{\alpha,K}^2
  \int_{\R^3}(1+|x|)^{-2K}\,dx
  <
  \infty.
\]
Only finitely many multi-indices occur below a fixed height \(m\), so \(u_0\in H^m_\sigma(\R^3)\) for every finite \(m\).
The case \(m=0\) says that the initial kinetic energy is finite.
\MajorThoughtBreak
\Needspace{18\baselineskip}
Now fix one such \(u_0\) and one viscosity \(\nu>0\).
Local theory produces a classical velocity--pressure history, and uniqueness joins every local piece into one maximal smooth interval \([0,T_*)\), where \(T_*\in(0,\infty]\).
\MajorThoughtBreak
The whole-space claim concerns this same history:
\begin{equation}\label{eq:WholeSpaceRegularityTarget}
  u,p\in C^\infty\!\bigl(\R^3\times[0,\infty)\bigr),
  \qquad
  \sup_{t\geq0}\int_{\R^3}|u(x,t)|^2\,dx<\infty.
\end{equation}
Energy is not the unresolved part.  \Cref{prop:ClassicalKineticEnergyIdentity} gives \(\norm{u(t)}_2\leq\norm{u_0}_2\) for every \(0\leq t<T_*\).
What remains is to prevent the smooth history from ending while that energy stays finite:
\[
  T_*=\infty.
\]
